\begin{theorem}
\label{theorem-ffdescent-projectivity}
Let $R \to S$ be a faithfully flat ring map.  Let $M$ be an $R$-module.
 If the $S$-module $M \otimes_R S$ is projective, then $M$ is projective.
\end{theorem}

\begin{proof}
We are going to construct a Kaplansky d\'evissage of $M$ to show that it is a
direct sum of projective modules and hence projective.  By Theorem
\ref{theorem-projective-direct-sum} we can write $M \otimes_R S =
\bigoplus_{i \in I} Q_i$ as a direct sum of countably generated $S$-modules
$Q_i$.  Choose a well-ordering on $M$.  Using transfinite recursion we are going
to define an increasing family of submodules $M_{\alpha}$ of $M$, one for each
ordinal $\alpha$, such that $M_{\alpha} \otimes_R S$ is a direct sum of some
subset of the $Q_i$.

\medskip\noindent
For $\alpha = 0$ let $M_0 = 0$.  If $\alpha$ is a limit ordinal and $M_{\beta}$
has been defined for all $\beta < \alpha$, then define $M_\alpha =
\bigcup_{\beta < \alpha} M_{\beta}$.  Since each $M_{\beta} \otimes_R S$ for
$\beta < \alpha$ is a direct sum of a subset of the $Q_i$, the same will be
true of $M_{\alpha} \otimes_R S$.  If $\alpha + 1$ is a successor ordinal and
$M_{\alpha}$ has been defined, then define $M_{\alpha + 1}$ as follows.  If
$M_{\alpha} = M$, then let $M_{\alpha +1} = M$.  Otherwise choose the smallest
$x \in M$ (with respect to the fixed well-ordering) such that $x \notin
M_{\alpha}$. Since $S$ is flat over $R$, $(M/M_{\alpha}) \otimes_R S = M
\otimes_R S/M_{\alpha} \otimes_R S$, so since $M_{\alpha} \otimes_R S$ is
a direct sum of some $Q_i$, the same is true of $(M/M_{\alpha}) \otimes_R
S$.   By Lemma \ref{lemma-adapted-submodule}, we can find a countably generated
$R$-submodule $P$ of $M/M_{\alpha}$ containing the image of $x$ in
$M/M_{\alpha}$ and such that $P \otimes_R S$ (which equals $\Im(P
\otimes_R S \to M \otimes_R S)$ since $S$ is flat over $R$) is a
direct sum of some $Q_i$. Since $M \otimes_R S = \bigoplus_{i \in I} Q_i$
is projective and projectivity passes to direct summands, $P \otimes_R S$ is
also projective.  Thus by Lemma \ref{lemma-ffdescent-countable-projectivity},
$P$ is projective.  Finally we define $M_{\alpha + 1}$ to be the preimage of $P$
in $M$, so that $M_{\alpha + 1}/M_{\alpha} = P$ is countably generated and
projective.  In particular $M_{\alpha}$ is a direct summand of $M_{\alpha + 1}$
since projectivity of $M_{\alpha + 1}/M_{\alpha}$ implies the sequence $0
\to M_{\alpha} \to M_{\alpha + 1} \to
M_{\alpha + 1}/M_{\alpha} \to 0$ splits.

\medskip\noindent
Transfinite induction on $M$ (using the fact that we constructed
$M_{\alpha + 1}$ to contain the smallest $x \in M$ not contained in
$M_{\alpha}$) shows that
each $x \in M$ is contained in some $M_{\alpha}$.  Thus, there is some large
enough ordinal $S$ satisfying: for each $x \in M$ there is $\alpha \in S$ such
that $x \in M_{\alpha}$.  This means $(M_{\alpha})_{\alpha \in S}$ satisfies
property (1) of a Kaplansky d\'evissage of $M$.  The other properties are clear
by construction.  We conclude
$M = \bigoplus_{\alpha + 1 \in S} M_{\alpha + 1}/M_{\alpha}$.
Since each $M_{\alpha + 1}/M_{\alpha}$ is projective
by construction, $M$ is projective.
\end{proof}
